% !TeX root = ../main.tex

\newlecture[][Sep 24, 2026]
\begin{theorem}[Independence and Expectation][5.7]
    If $X$ and $Y$ are independent with $EX, EY < \infty$, then 
    \[E(XY) = (EX)(EY)\]
    Proof: 
    \begin{quote}
        First, consider both $X$ and $Y$ are simple, so that
        \[X= \sum_{i=1}^{n} x_i I(A_i) \quad \text{and} \quad Y= \sum_{j=1}^{m} y_j I(B_j)\]
        \begin{align*}
            E(XY) &= E[\sum_{i=1}^{n}\sum_{j=1}^{m} x_i y_j \cdot [\left\{A_i \bigcap B_j\right\}]\\
            &\phantom{=}\ \sum_{i=1}^{n}\sum_{j=1}^{m} x_i y_j P(x_i \in A_i, Y_j \in B_j)]\\
            &=\sum_{i=1}^{n}\sum_{j=1}^{m} x_i y_i P(\left\{\omega: X(\omega) \in A_i, Y(\omega) \in B_j\right\})\\
            &=\sum_{i=1}^{n}\sum_{j=1}^{m} x_i y_j P(x \in A_i)\cdot P(Y\in B_j)\\
            &=\underbrace{\sum_{i=1}^{n}x_i\cdot P(x\in A_i)}_{EX}  \underbrace{\sum_{j=1}^{m}y_j\cdot P(Y \in B_j)}_{EY}\\
            &= (EX)(EY)
        \end{align*}
    \end{quote}
\end{theorem}

Next, consider the case that $X$ and $Y$ are non-simple and non-negative, then we let $Z_1\geq X, Z_2 \geq Y$, where $Z_1$ and $Z_2$ are simple. Notice that
\[Z_1 \geq X, Z_2 \geq Y \text{ implies that } Z_1 \cdot Z_2 \geq XY\]
\begin{align*}
    (EX)(EY) &= \inf \left\{EZ_1: Z_1 \text{ simple, } Z_1 \geq X \text{ a.s.}\right\}\\
    &\mathrel{\phantom{=}}\inf\left\{EZ_2: Z_2\text{ simple, } Z_2 \geq Y \text{ a.s.}\right\}
\end{align*}
\[ \inf\left\{(EZ_1) \cdot(EZ_2)\right\} \quad \text{because of indepegndence and $Z_1, Z_2$ simple}\]
\[\geq \inf\left\{EZ_1Z_2, Z_1 \text{ and }Z_2 \text{ simple and indep, } Z_1 Z_2\geq XY \text{ a.s.}\right\}\]

Additionally, if $Z_1$ and $Z_2$ are simple and 
\[Z_1 \leq X \text{ a.s.} \qquad Z_2 \leq Y \text{ a.s.}\]
Similarly, $(EX)(EY = \sup\left\{EZ_1Z_2, Z_1 \text{ and }Z_2 \text{ are simple and indep, }Z_1 \leq X, Z_2 \leq Y \text{ a.s.}\right\})$


Previously, we proved that
\[\inf\left\{EZ_1Z_2, \quad Z_1, Z_2 \text{ are simple and indep with }Z_1 \geq X, Z_2 \geq Y \text{ a.s. }\right\} = E(XY)\]
and
\[\sup\left\{EZ_1Z_2, \quad Z_1, Z_2 \text{ are simple and indep with }Z_1 \leq X, Z_2 \leq Y \text{ a.s. }\right\} = E(XY)\]


Finally, for $X$ and $Y$ that are integrable, first realize that if $X$ and $Y$ are independent, then $X^+, X^- Y^+, Y^-$ are also independent.
\[X = X^+ - X^- \qquad Y = Y^+ - Y^-\]
\begin{align*}
    E(XY) &= E[(X^+ - X^-)(Y^+ - Y^-)]\\
    &=E[X^+Y^+ - X^+ Y^- - X^- Y^+ + X^- Y^-]\\
    &= E(X^+Y^+) - E(X^+Y^-) - E(X^- Y^+) + E(X^- Y^-)\\
    &= E(X^+)E(Y^+) - E(X^+)E(Y^-) - E(X^-)E(Y^+) + E(X^-)E(Y^-)\\
    &=(EX)(EY)
\end{align*}

\begin{definition}[Variance and Covariance][5.8]
    The variance of a random variable $X$ is defined as
    \[\Var(X) = E[(X-EX)^2]\]
    The covariance of random variables $X$ and $Y$ is 
    \[\Cov(X,Y) = E[(X-EX)(Y-EY)]\]
    The correlation coefficient is
    \[\text{Corr}(X,Y) = \frac{\Cov(X,Y)}{\sqrt{\Var(X)\Var(Y)}}\]
\end{definition}
